\section{Methodology}
\label{sec:methodology}

We adapt the standard \nih evaluation paradigm~\citep{liu2023lost} from factual retrieval to adversarial injection, embedding prompt injection attacks at controlled positions within documents of varying lengths and measuring whether the model follows the injected instruction.

\subsection{Task and Threat Model}
\label{sec:task}

We consider a summarization task where the model receives a system prompt (``You are a helpful assistant that summarizes documents'') and a user prompt containing an external document to summarize.
The attacker controls a portion of the document and embeds an adversarial instruction, but cannot modify the system prompt or user instruction.
This threat model reflects real-world scenarios such as RAG pipelines, email summarization, and document processing, where the LLM ingests untrusted content.

\subsection{Document Construction}
\label{sec:documents}

\para{Filler text.}
We curate 16 coherent paragraphs covering science, history, technology, and nature topics in an encyclopedia-like style.
For each experimental condition, paragraphs are sampled (with seed 42) and concatenated to reach the target document length.
This ensures realistic, coherent context rather than random noise.

\para{Injection placement.}
Given a target document length $L$ and injection position $p \in [0, 1]$, we split the filler text at position $\lfloor p \cdot L \rfloor$ tokens and insert the adversarial payload.
Position $p=0$ places the injection at the document start, $p=0.5$ at the middle, and $p=1.0$ at the end.

\para{Experimental grid.}
We test 7 document lengths $L \in \{500, 1000, 2000, 4000, 8000, 16000, 32000\}$ tokens and 7 injection positions $p \in \{0.0, 0.1, 0.25, 0.5, 0.75, 0.9, 1.0\}$, yielding 49 length-position combinations per attack.

\subsection{Attack Design}
\label{sec:attacks}

We design 10 adversarial prompts in two categories (\secref{sec:attack_appendix}):

\para{Direct attacks} (5 attacks) attempt explicit instruction override.
These include instruction overrides (``Ignore all previous instructions\ldots''), role manipulation (``[ADMIN MODE ACTIVATED]\ldots''), context confusion using document-end delimiters, social engineering via flattery, and developer-note instruction leaks.

\para{Subtle attacks} (5 attacks) blend into the document's expected conventions.
These include fake editor notes, citation requirements, formatting guidelines, compliance notices, and author preference statements.
Each subtle attack frames its payload as legitimate document metadata rather than as an instruction override.

Each attack has a unique success pattern---a specific string (\eg \texttt{INJECTION\_SUCCESSFUL}, \texttt{CATEGORY-BLUE}) that appears in the model's output only if the injection succeeds.

\subsection{Models and Configuration}
\label{sec:models}

We test two OpenAI models: \gptmini (full factorial grid: 7 lengths $\times$ 7 positions $\times$ 10 attacks $\times$ 2 trials = 980 experiments) and \gptfull (reduced grid: 5 lengths $\times$ 5 positions $\times$ 10 attacks $\times$ 1 trial = 250 experiments).
All experiments use temperature 0.0 for deterministic outputs and a maximum of 300 output tokens.

\subsection{Evaluation}
\label{sec:evaluation}

\para{Attack Success Rate (ASR).}
We evaluate injection success via pattern matching: if the model's response contains the attack's success pattern (case-insensitive), the attack is scored as successful.
This approach is conservative---it may miss partial compliance but ensures high precision.
We manually verified 20 randomly sampled results (10 successes, 10 failures) and found 100\% agreement with the automated scoring.

\para{Statistical tests.}
We use the Kruskal-Wallis $H$-test for omnibus comparisons across length and position levels, the Mann-Whitney $U$-test for pairwise comparisons between attack categories, and Spearman rank correlation for monotonic trends.
All tests use $\alpha = 0.05$.
